\begin{topic}[id=logistics_drones,
             difficulty={Mittel},
             type={Anwendungsanalyse + Bewertung},
             prerequisites={Grundlagen eingebetteter Systeme; Interesse an autonomen Transport- und Assistenzsystemen},
             practice={optional},
             owner={Dozententeam},
             updated={2026-04-09},
             tags={ Drohnen, Logistik, BVLOS, Sicherheit, Navigation, Regulierung }]{Logistikdrohnen}

\TopicDescription{Drohnen versprechen schnelle, flexible Logistik – von Werkslogistik bis Zustellung. Wir beleuchten Technik (Navigation, Energie, Redundanz), Betriebskonzepte (VLOS/BVLOS) und Sicherheits-/Regulierungsaspekte. Ziel: realistische Bewertung der Einsatzfähigkeit und Systemarchitekturen.}

\TopicLeitfragen{\item Welche Einsatzszenarien sind wirtschaftlich tragfähig?
\item Wie adressiert man BVLOS-Risiken?
\item Welche Sensorik/Redundanz ist sinnvoll?
}
\TopicScope{\item Keine Rechts-/Genehmigungsdetails im Land-spezifischen.
\item Keine Flugsteuerungs-Programmierung.
}
\TopicPractice{Optionale Praxisidee: Missionsarchitektur und Risikoanalyse für ein Campus-Szenario.}

\TopicKeywords{Drohnen, Logistik, BVLOS, Sicherheit, Navigation, Regulierung}

\nocite{easaDrones,droneSurvey}
\end{topic}
